\section{The direct proof-calculus source}
\label{sec:source}

This section records the inherited source mechanism
\cite{shared,flow} so that the consistency interpretation
is explicit. The new compression acts on the whole
source, not merely a local arithmetic relation.

The fixed numerical calculus $C_{222}$ uses implication
code $I(a,b)=2(a+b)(a+b+1)+4b$, with offsets $1,2,3$
for membership, universal quantification and equality.
Negation is $I(a,1)$. Appendix~\ref{app:calculus}
gives all constructors and twelve rule branches.
Truth is the freely available code $0$; contradiction
is code $1$. The numerical calculus is calibrated
to syntactic ZFC refutability by the named inherited
theorems in Section~\ref{sec:verification}.

\subsection{Finite flow certificates}

A padded rule with premises $p,q$ and conclusion $v$
has vector $e_v-e_p-e_q$, where $e_0=0$ and positive
labels are independent basis vectors. A finite list
of valid instances satisfying
\begin{equation}
 \sum_i(e_{v_i}-e_{p_i}-e_{q_i})=e_1
\label{eq:flow}
\end{equation}
already implies derivability of $1$.
Let $D$ be the closure under the listed rules of
truth and their axiom conclusions. If $1\notin D$,
sum the coordinates outside $D$. A conclusion
outside $D$ has at least one premise outside $D$,
so every rule contributes at most zero there.
The right side contributes one, a contradiction.
Only finitely many labels occur, so this is a
finite sum. Cycles and repeated premises cause
no exception.

Conversely a finite proof tree gives the flow by
counting its nodes. Expand it starting from the
open obligation $1$, replacing each conclusion
by its selected premises.

\subsection{A carrier recovering the exact flow}

Let $t,m\ge1$ be the number of instances and a
bound on their labels, and put $M=(m+1)^t$.
Set
\[
 A(X)=\prod_{i<t}(1+v_iX),\qquad
 C(X)=(1+X)\prod_{i<t}(1+p_iX)(1+q_iX).
\]
Both have nonnegative coefficients and constant
term one. Then
\begin{equation}
 A(M)=C(M)\quad\Longleftrightarrow\quad A(X)=C(X).
\label{eq:carrier}
\end{equation}
If $\deg C\ge t+1$, its evaluation is at least
$M^{t+1}$, whereas
\[
 A(M)\le(1+mM)^t<((m+1)M)^t=M^{t+1}.
\]
Thus each product has at most $t$ nonconstant
factors and coefficient sum at most $M$.
Every individual coefficient is strictly below
$M$, because the constant term is one and
a nonconstant product has another positive
coefficient. Every coefficient of $A-C$
therefore has absolute value below $M$.
Its least nonzero coefficient, if any, would
be divisible by $M$, by reducing its zero
evaluation after removing the corresponding
power of $M$. This is impossible.

Unique factorization in $\mathbb Q[X]$ now
recovers the exact multiplicities of the
distinct factors $1+jX$ for $j>0$, giving
\eqref{eq:flow}. No coprimality of their
integer values is assumed.

Natural accumulator states satisfy
\begin{equation}
 S_0=M+1,\quad S_t=1,\quad
 S_{i+1}(1+v_iM)=S_i(1+p_iM)(1+q_iM).
\label{eq:accumulator}
\end{equation}
Positive factors force every state positive,
and telescoping gives \eqref{eq:carrier}.
For completeness take $S_i$ to be the product
of the weights of current open obligations
in the proof-tree expansion. Every update
is then an integer identity.

One beta code represents the accumulator:
$\beta(x,y,i)=x\bmod(1+(i+1)y)$.
Choose $y$ sufficiently large and divisible
by every relevant positive index difference.
The moduli are pairwise coprime and exceed
all sequence entries, so CRT supplies $x$.
The choice order is finite proof, $t,m$,
carrier, accumulator, then $x,y$.
The later witness bound is independent of
the earlier label bound $m$.

\subsection{Thirteen local coordinates}

The local row uses three labels $p,q,v$,
four schema coordinates $a,b,c,d$ and six
beta/accumulator coordinates $S,S',u,r,u',r'$.
For $D_i=1+(i+1)y$ its arithmetic rows are
\begin{align*}
 x&=S+D_i u,&D_i&=S+r+1,\\
 x&=S'+(D_i+y)u',&D_i+y&=S'+r'+1,\\
 S'(1+vM)&=S(1+pM)(1+qM).
\end{align*}
The first four rows give exact consecutive
beta digits, since the divisors are positive.

For each rule branch square and sum its
residuals; multiply the twelve branch sums
to impose their disjunction. Add the
squares of the five other local rows.
This is a fixed nonnegative polynomial
$L(i,\mathbf z,\mathbf w)$.
Its implementation uses nonnegative-coefficient
circuits $E,F$ with $L=E-F$. For multiplication,
\[
 (E_1,F_1)(E_2,F_2)=
 (E_1E_2+F_1F_2,\ E_1F_2+F_1E_2).
\]
Addition and subtraction use componentwise
addition and the appropriate swap. The
difference is preserved and both components
are monotone in natural inputs.

\section{Removing the bounded universal quantifier}
\label{sec:crt}

Give the local witnesses bounds $u_j$ and
choose $V\ge u_j$ for every $j$. Define
\[
 H=E(t,\mathbf z,\mathbf u)+F(t,\mathbf z,\mathbf u)+t+V+1,
 \qquad T=H!,\qquad
 B_0=\prod_{i<t}(1+(i+1)T).
\]
Existence of bounded witnesses for every
index $i<t$ is equivalent to natural
$C,A_1,\ldots,A_d$ satisfying
\begin{align}
 B_0&\mid\binom{A_j}{u_j+1}
             &&\text{for each }j,\label{eq:crt-range}\\
 B_0&\mid1+(C+1)T,\label{eq:crt-index}\\
 E(C,\mathbf z,\mathbf A)&\equiv
       F(C,\mathbf z,\mathbf A)\pmod{B_0}.
       \label{eq:crt-row}
\end{align}

\begin{proof}[Soundness]
Fix $i<t$ and a prime divisor $p$ of
$D_i=1+(i+1)T>1$.
It is coprime to $T=H!$, so $p>H$.
The integer identity
\[
 (u_j+1)!\binom{A_j}{u_j+1}
       =\prod_{r=0}^{u_j}(A_j-r)
\]
also holds for $A_j\le u_j$, when both
sides vanish. Primality supplies a
$w_j\in[0,u_j]$ with $A_j\equiv w_j\pmod p$.
Comparing the index divisibility with
$D_i\equiv0$ and cancelling $T$ gives
$C\equiv i\pmod p$.
The row congruence yields
$E(i,\mathbf z,\mathbf w)\equiv
F(i,\mathbf z,\mathbf w)\pmod p$.
Monotonicity places both natural values
strictly below $p$, hence they are equal.
\end{proof}

\begin{proof}[Completeness]
A common divisor of $D_i,D_j$ for $i<j<t$
divides $(j-i)T$ and is coprime to $T$,
so it divides $j-i$. But $j-i\mid H!$,
making the common divisor one.
CRT now chooses $C\equiv i$ and
$A_j\equiv w_j^{(i)}\pmod{D_i}$
for valid bounded row witnesses.
Their falling-factorial products vanish
modulo every $D_i$. Since
$(u_j+1)!\mid T$ and $\gcd(B_0,T)=1$,
that factorial cancels, giving the
binomial divisibility. The index and
polynomial congruences follow by
substitution and pairwise coprimality.
\end{proof}

Soundness deliberately uses a prime
divisor of each modulus. A product
vanishing modulo a composite number
does not by itself supply a small root.
Completeness deliberately establishes
$H\ge u_j+1$ before factorial cancellation.

For the actual source, the three labels
have bound $m$ and the other ten local
coordinates have bound $V=m+U_0$ for an
extra natural $U_0$. There are thirteen
range-binomial graphs. Endpoint equations
directly encode $S_0=M+1,S_t=1$.
Writing $t=t_0+1,m=m_0+1$ loses no proof:
zero length contradicts the endpoints,
and zero label bound makes every update
constant. Encoded coordinates may be
shifted to $A_j=\alpha_j+1$, since adding
a multiple of $B_0$ preserves the
falling-factorial divisibility and its
factorial cancellation.

\subsection{Eliminating the remaining arithmetic macros}

The explanatory product is replaced by
\[
 Q_0=(1+(t+1)T)^t+2,\qquad
 B_0=
 \left(T^t\,t!\binom{Q_0+t}{t}\right)\bmod(TQ_0-1).
\]
The numerator is
$\prod_{i=1}^t(TQ_0+iT)$ and reduces
to the intended product modulo $TQ_0-1$.
That product is smaller than the
positive modulus, so it is the exact
remainder.

For a natural $n$, let
$\rho=(n+2)^{n+2}$,
$C=\binom{\rho+n}{n}$ and
$P=(\rho+1)^n$. Then
\[
 n!=\left\lceil P/C\right\rceil.
\]
For $n\ge2$,
\[
 \frac PC=n!\prod_{i=1}^n
 \left(1-\frac{i-1}{\rho+i}\right),
 \qquad
 0\le n!-\frac PC
 <\frac{n!\,n(n-1)}{2(\rho+1)}<1.
\]
The cases $0,1$ are immediate.
Natural $r,s$ impose the ceiling by
$P+C=Cf+r+1$ and $C=r+s+1$.
Each other division uses the analogous
quotient/remainder equations.
All binomials use the shared radix,
and all powers use the shared Pell
certificate. No arithmetic macro
remains in the final polynomial.

\subsection{The positivity order at arbitrary zeros}

All twenty-three exponent groups are
positive from structural shifts and
nonnegative arithmetic, before any
power output is known correct.
Consequently every group has
$Y_j\ge e_j>0$ and the common index
certificate recovers all its Pell
pairs first.

Twenty-seven of the twenty-eight
power bases are structurally positive.
The remaining base is the factorial
output $T$, which must be proved
positive rather than silently assumed.
Apply the local power argument first
to the structurally positive bases.
In particular the factorial numerator
$P=(\rho+1)^H$ is positive.
The retained ceiling rows
$P+C=CT+r+1$ and $C=r+s+1$
would imply $P\le0$ if $T=0$.
Hence $T>0$ and the final power
$T^t$ is valid.

The common radix power has literal
base two and a positive exponent, so
its correctness is available before
the binomial digit argument.
Its coefficient bound depends only
on the assigned natural upper
arguments, not on their prior
semantic correctness.
This order resolves the apparent
cycle introduced by deferring
factorial and binomial outputs.
